\section{Fourteen first-order shells and all their periodic lifts}
Let $\mathcal S\subset(0,1)$ be the union of the fourteen open intervals
in the first column of the table below.  For every prime $p>29$ with
$p^2>30n$ and $\{n/p\}\in\mathcal S$, one has
\[
 v_p(Q_n)\ge1,\qquad v_p(P_{3,n})\ge-1,\qquad
 v_p(P_{4,n})\ge-2.                                    \tag{PS}
\]
Thus the extra prime divisor may include every such prime, not only those
in the original shell $8n<p<11n$.  These are first-order divisibility
statements.  No unproved second-order enhancement is included.

\paragraph{The finite root-count certificate.}
Put $n/p=t+f$ with $t\in\mathbb Z_{\ge0}$ and $0<f<1$, and write
$y=k/p$ for a residue $0\le k<p$.  At $t=0$, the three numerator high-root
intervals have integer endpoints
\[
 [\lfloor uf-y\rfloor+1,\ \lfloor(u+10)f-y\rfloor],\quad u=0,1,2,
\]
and the two denominator intervals have endpoints
\[
 [\lceil-17f-y\rceil,\ \lfloor-f-y\rfloor],\qquad
 [\lceil-16f-y\rceil,\ \lfloor-2f-y\rfloor].
\]
Let $h$ be the total numerator count minus the total denominator count,
and put $\epsilon_0=\lfloor2f-y\rfloor$.
For every $f\in\mathcal S$ away from Farey walls of order $29$ (in particular, every admissible $f=\{n/p\}$), the following two assertions hold:
$h\ge-1$ at every residue; and wherever $h=-1$, subtracting $\epsilon_0$
from every one of the ten endpoints yields the single row of the table
corresponding to $f$.  Notation $[L,U]$ includes every integer from $L$ to
$U$ with multiplicity; overlapping intervals retain their multiplicities.
\[
\begin{array}{c|ccc|cc}
 f\text{ interval}&N_0&N_1&N_2&D_0&D_1\\\hline
 (1/11,1/8)&[1,1]&[1,1]&[1,1]&[-1,0]&[-1,0]\\
 (3/13,4/17)&[0,1]&[0,2]&[1,2]&[-4,-1]&[-4,-1]\\
 (3/11,5/18)&[0,2]&[0,2]&[1,2]&[-5,-1]&[-4,-1]\\
 (4/13,5/16)&[0,2]&[1,3]&[1,3]&[-5,-1]&[-5,-1]\\
 (4/11,7/19)&[0,3]&[1,3]&[1,4]&[-6,-1]&[-6,-1]\\
 (5/13,7/18)&[1,3]&[1,4]&[1,4]&[-6,-1]&[-6,-1]\\
 (5/12,8/19)&[0,3]&[1,4]&[1,4]&[-7,-1]&[-7,-2]\\
 (5/11,8/17)&[0,3]&[0,4]&[1,4]&[-8,-2]&[-8,-2]\\
 (6/11,9/16)&[0,4]&[0,5]&[1,5]&[-10,-2]&[-9,-2]\\
 (7/12,10/17)&[-1,4]&[0,5]&[1,5]&[-11,-2]&[-10,-3]\\
 (7/11,11/17)&[-1,5]&[0,5]&[1,6]&[-12,-2]&[-11,-3]\\
 (7/10,12/17)&[-1,5]&[0,6]&[1,7]&[-13,-2]&[-12,-3]\\
 (8/11,11/15)&[-1,5]&[0,6]&[1,7]&[-13,-3]&[-13,-3]\\
 (10/11,16/17)&[-1,7]&[0,8]&[1,9]&[-17,-3]&[-16,-4]
\end{array}
\]
This finite table can be verified entirely with rational arithmetic.
All changes of an endpoint occur on a line
$y=mf+c$, with $m\in\{0,1,2,10,11,12,-1,-2,-16,-17\}$ and integer $c$.
The crossing abscissas in $(0,1)$ therefore have reduced denominator at
most $29$.  Partition the $f$-axis by the Farey sequence of order $29$;
there are $270$ open cells.  At the midpoint of each, test every
$y$-boundary and every intervening open interval, including $y=0$ and
excluding $y=1$.  The root data are constant on each tested cell, so this
finite verification proves the table, rather than testing selected
values of $n$.  The supplied certificate checks all $5400$ such cells
and boundaries; $47$ of the Farey cells form the fourteen displayed
intervals.  A missing Farey boundary cannot occur for $f=\{n/p\}$ when
$p>29$; its reduced denominator would have to be $p$.

\paragraph{Why the same table proves every lift.}
Changing $f$ to $t+f$ adds $30t$ numerator roots and $30t$ denominator
roots at every residue.  Hence $h$ and its minimal support are unchanged.
Also $\epsilon=\lfloor2(t+f)-y\rfloor=2t+\epsilon_0$.
For a table entry $N_u=[L_u,U_u]$, the high-root interval after subtracting
$\epsilon$ is
\[
 [(u-2)t+L_u,\ (u+8)t+U_u].
\]
For $D_0=[L_3,U_3]$ and $D_1=[L_4,U_4]$ it is respectively
\[
 [-19t+L_3,-3t+U_3],\qquad[-18t+L_4,-4t+U_4].
\]
These five intervals depend only on $t$ and the table row, and are
independent of the residue $k$ on the entire minimal support.  Let
$U_t(V)$ be the monic rational function having these numerator and
denominator roots.  It has degree $-1$.

As in the broad-shell proof,
$\mathcal Q=(X^p-X)\overline{\mathcal N}/\overline{\mathcal D}$ is a
polynomial of degree $p-2$.  At a minimal-support residue the local
factorization is
\[
 pF_n\bigl(k+p(V+\epsilon)\bigr)
       \equiv-\mathcal Q(k)U_t(V)\pmod p;
\]
outside that support its reduction is zero.  Distinct high roots differ
by at most $\lfloor29n/p\rfloor<p$, so all local partial-fraction
operations have $p$-integral coefficients.  A denominator root $l$ of
$U_t$ has harmonic index $-l$ after extracting the multiples of $p$ from
$H_{j+2n}^{(r)}$.  These indices are at most $19n/p<p$, so their harmonic
numbers are $p$-integral too.  The resulting leading local readout is
therefore a constant depending on $t$ and the table row, multiplied by
$\mathcal Q(k)$.  The sum over $k$ vanishes since
$\deg\mathcal Q<p-1$.  This proves (PS) for every periodic lift.

The normalization valuation
$c(f)=2\lfloor15f\rfloor-3\lfloor10f\rfloor$ is unchanged by the lift.
On the table intervals the baseline profile is exactly $c(f)+1$ for both
readouts, including every lifted interval; (PS) raises it to $c(f)+2$.
The same rational certificate checks these equalities against the
baseline chamber table.

\paragraph{Arithmetic density.}
If
\[
 H_n^*=\prod_{\substack{p>29,\ p^2>30n\\\{n/p\}\in\mathcal S}}p,
\]
then the prime number theorem gives
\[
 \lim_{n\to\infty}\frac{\log H_n^*}{n}
 =\sum_{(a,b)\subset\mathcal S}\sum_{t=0}^{\infty}
      \left(\frac1{t+a}-\frac1{t+b}\right).
\]
The discarded primes have logarithmic total $o(n)$.  All terms are
positive.  At $t=0$ the exact sum is $195197/55440$; summing only
$t=0,\ldots,999$ gives the exact rational lower bound recorded by the
certificate, and in particular
\[
 \liminf_{n\to\infty}\frac{\log H_n^*}{n}
       >3.6747281958.
\]
Thus the previously certified baseline density $26.2366642995$ improves
to a certified density greater than $29.9113924953$.
